\section{Rancangan Eksperimen}
\subsection{Korpus dan Pertanyaan}
Korpus \tech{User Manual} terdiri atas sembilan manual Netgear dengan total 1.400 halaman. Model produk yang dicakup adalah CM2000, LM1200, R6350, R7000, RAX120, RAX50, RAXE500, RBK852, dan XR500. Dari struktur bagian dan \tech{block} teks dipilih 100 \tech{intent} menggunakan \tech{seed} 20260831 dengan mempertimbangkan keterwakilan dokumen dan posisi halaman. Setiap \tech{intent} dirumuskan menjadi satu pertanyaan bahasa Inggris dan satu pertanyaan bahasa Indonesia, menghasilkan 200 pertanyaan.

Setiap pasangan bahasa menggunakan satu \tech{block} sumber yang sama sebagai \tech{context reference}. Pertanyaan dan jawaban acuan disusun dengan bantuan model berdasarkan paket konteks. Karena dokumen sumber berbahasa Inggris, pertanyaan bahasa Indonesia mengukur pencarian lintas bahasa. Kedua bahasa diberi bobot yang sama, tetapi 200 pertanyaan tersebut mewakili 100 kebutuhan informasi berpasangan.

Korpus peraturan dibatasi pada bidang pendidikan di Indonesia. Dari populasi awal 8.745 dokumen, 4.459 memiliki keluaran \tech{parser} yang siap diproses secara vektor. Pemeriksaan pasangan sumber dan \tech{hash} menghasilkan 4.453 dokumen valid sebagai kerangka pemilihan. Dokumen yang terhubung melalui relasi perubahan atau pencabutan dikelompokkan menjadi \tech{primary sampling unit} (PSU).

Strata utama dibentuk berdasarkan jenis relasi dan komposisi status hukum. Bentuk peraturan, tahun, dan panjang dokumen digunakan untuk memeriksa keragaman. Pemilihan PSU dilakukan dengan \tech{probability proportional to size} di dalam strata. Seluruh anggota PSU terpilih dipertahankan agar hubungan antarperaturan tidak terputus. Hasilnya adalah 255 dokumen dari 217 PSU.

Sebanyak 25 PSU menyediakan masing-masing empat pertanyaan, sehingga terdapat 100 pertanyaan hukum. Model menyusun pertanyaan dan jawaban acuan dari konteks yang disiapkan, kemudian \tech{context reference} ditetapkan pada satu atau beberapa ID ayat kanonis. Sebanyak 30 pertanyaan juga memiliki referensi hubungan perubahan, pencabutan, atau garis keturunan untuk evaluasi relasi. Korpus dan pertanyaan dibekukan sebelum dibandingkan antarstrategi.

\subsection{Kondisi Pengujian}
Eksperimen bersifat \tech{retrieval-only}. Pembentukan jawaban dan \tech{reranking} dinonaktifkan. Konfigurasi API dan \tech{worker} dipertahankan. Perbandingan menggunakan pendekatan \tech{one factor at a time} (OFAT), yaitu mengganti satu komponen sementara komponen pembanding dibuat tetap.

Pada \tech{User Manual}, perbandingan \tech{chunker} menggunakan \tech{hybrid indexer} dan BAAI/BGE-M3. Strategi \tech{structure-aware} memakai \tech{target text block} dengan \tech{parent context} dan batas 512 karakter. Strategi \tech{recursive} dan \tech{sliding-window} menggunakan 512 karakter dengan \tech{overlap} 51 karakter. Perbandingan \tech{indexer} mempertahankan \tech{structure-aware chunker}, lalu mengganti strategi menjadi \tech{sparse}, \tech{dense}, atau \tech{hybrid}.

Pada peraturan, kondisi acuan menggunakan \tech{legal structure-aware chunker}, BAAI/BGE-M3, dan \tech{dense indexer}. Penggantian model ke Indo-LegalBERT, penggantian \tech{chunker}, dan penggantian \tech{indexer} dilakukan secara terpisah. Pengisian konteks induk saat pengukuran dinonaktifkan. Hal ini perlu dibedakan dari konteks struktural yang telah dibawa ketika \tech{chunk} dibentuk.

\subsection{Ukuran Evaluasi}
Ukuran utama adalah \tech{recall@K}, yaitu proporsi \tech{context reference} yang ditemukan dalam $K$ hasil teratas. Untuk satu pertanyaan $q$, definisinya dapat dituliskan sebagai
\begin{equation}
R_q@K=\frac{|G_q\cap C_q(K)|}{|G_q|},
\end{equation}
dengan $G_q$ sebagai himpunan referensi acuan dan $C_q(K)$ sebagai referensi yang dicakup hasil teratas. Pencocokan dilakukan melalui identitas sumber, bukan kesamaan teks jawaban. Pada \tech{User Manual} yang memiliki satu referensi per pertanyaan, ukuran ini menunjukkan keberhasilan menemukan referensi tersebut. Pada peraturan, satu pertanyaan dapat membutuhkan beberapa referensi.

\tech{Precision} dan \tech{nDCG} dilaporkan sebagai diagnostik kepadatan hasil relevan dan urutan pencarian, sedangkan pemilihan konfigurasi mengutamakan \tech{recall}. \tech{KG recall} dibaca terpisah pada 30 pertanyaan relasi. Ukuran tersebut menilai penemuan sumber yang memuat hubungan target dan tidak sama dengan evaluasi menyeluruh ketepatan ekstraksi graf.

Nilai $K$ diperiksa pada rentang 5 hingga 60 dengan selang lima. Pada peraturan, titik operasional dipilih ketika kenaikan \tech{recall} paling besar 0,02 untuk setiap penambahan lima kandidat dan tetap demikian pada seluruh nilai berikutnya. Perbandingan model memakai $K=45$, sedangkan perbandingan \tech{chunker} dan \tech{indexer} memakai $K=25$. Pada \tech{User Manual}, titik operasional \tech{dense} berada pada $K=40$, sementara perbandingan keseluruhan memakai $K=45$ yang sama bagi seluruh kondisi.
